\documentclass[a4paper,12pt]{report}
\usepackage[utf8]{inputenc} %s ovim mogu pisati hrv znakove direktno na linux i win. platformama (ali samo u utf-8 encodingu)
\usepackage{amsmath}
\usepackage[croatian]{babel}
\usepackage{xspace}
\usepackage{url}
\usepackage{hangcaption}
\usepackage[croatian,refpage]{nomencl}
\makenomenclature
\usepackage[diplomski,avant]{fsbthesis}
    % ili npr. \usepackage[avant,draft]{fsbthesis} %za doktorat - draft
	%   ispisuje broj verzije teksta (promijeni naredbu \ver) i datum
    % ili npr. \usepackage[zavrsni,avant]{fsbthesis} %za zavrsni rad
    % ili npr. \usepackage[diplomski,avant]{fsbthesis} %za diplomski rad
    % ili npr. \usepackage[magistarski,avant]{fsbthesis} %za mag.rad
    % ili npr. \usepackage[klasik]{fsbthesis} %za klasican izgled doktorata
    % ili npr. \usepackage[magistarski,klasik]{fsbthesis} %za klas.mag.rad
	% ... za opis svih opicija pogledaj fsbthesis.sty
%\usepackage{t1enc}

\vfuzz=2.7pt % Don't bother to report overfull boxes < 2.7pt
%\graphicspath{{slike/}} % dodaj poseban direktoriji za slike (vrlo prakticno)

\title{NASLOV DIPLOMSKOG RADA}  %Declares the document's title.
\author{ime prezime}      % Declares the author's name.
\date{Zagreb, svibanj 2003.}      % Deleting this command produces today's date.

\makeatletter
\renewcommand{\FSBautor}{Moje ime i prezime}
\renewcommand{\FSBmentor}{prof.~dr.~sc.~ime prezime} %ili
%	\begin{minipage}[t]{7cm}
%		dr.~sc.~ime prezime,\\
%		prof.~u mirovini
%	\end{minipage}}
\renewcommand{\FSBbibmentor}{
		dr.~sc.~ime prezime,
		prof. u mirovini}
\renewcommand{\FSBgodina}{2003.}
\renewcommand{\FSBbibudk}{xxx.xx}
\renewcommand{\FSBkljucnerijeci}{popis klju\v{c}nih rije\v{c}i: maksimalno do deset \ldots}
\renewcommand{\FSBkeywords}{list of the keywords: up to ten words}
\renewcommand{\FSBbibznpolje}{zrakoplovstvo, raketna i svemirska tehnika}
\renewcommand{\FSBbibdatumobr}{1. travnja 2005.}
\renewcommand{\FSBbibpovjerenstvo}{prvi \v{c}lan povjerenstva, \newline
                                   drugi \v{c}lan povjerenstva, \newline
                                   tre\'{c}i \v{c}lan povjerenstva, \newline
                                   \v{c}etvrti \v{c}lan povjerenstva, \newline
                                   peti \v{c}lan povjerenstva }
\renewcommand{\FSBpdfbrojstrana}{111}
\renewcommand{\FSBpdfbrojslika}{111}
\renewcommand{\FSBpdfbrojtablica}{111}
\renewcommand{\FSBpdfbrojrefs}{111}

\renewcommand{\ver}{\mbox{ver.2-01}}

\renewcommand{\thestrana}{11}


% komande za popis oznaka u nomenklaturnom stilu 
\newcommand{\dottfill}{%
  \leavevmode
  \cleaders \hb@xt@ .8em{\hss.\hss}\hfill
  \kern\z@}
\ifpdf
	\renewcommand{\pagedeclaration}[1]{\dottfill \hyperpage{#1}}%
\else
\renewcommand{\pagedeclaration}[1]{\dottfill \nobreakspace#1}%
\fi
\renewcommand{\nomname}{Popis oznaka}
\setlength{\nomitemsep}{-\parsep}
\addtolength{\nomitemsep}{0.2pt}
\renewcommand{\nomgroup}[1]{%
	\ifthenelse{\equal{#1}{K}}{\bigskip\par\item[\textbf{\large Kratice}]}{%
	\ifthenelse{\equal{#1}{I}}{\bigskip\par\item[\textbf{\large	Indeksi}]}{%
	\ifthenelse{\equal{#1}{T}}{\bigskip\par\item[\textbf{\large	Akcenti}]}{}}} }
%\newcommand{\nomunit}[1]{%
%	\renewcommand{\nomentryend}{\hspace*{\fill}[#1]}
\makeatother

% osobne definicije
\def\AR{\hbox{\textit{A}\hskip-.38em \rm R}} % aspect ratio sign
\def\Vinf{\hbox{$V_{\infty}$}} %brzina strujanja


\begin{document}
%
\pagenumbering{roman} % na pocetku rada stranice brojim rimski od (i)
%
\fsbkorice
%
\fsbprvilist
%
%\setcounter{brojstrana}{11}
\fsbbiblist
%
\begin{zahvala}
	\emph{To Donald E. Knuth \dots i druge zahvale \dots}
\end{zahvala}

\begin{izjava}
	Izjavljujem da sam ovaj rad radio samostalno koristeći stečena
	znanja tijekom studija i navedenu literaturu.
\end{izjava}

%-- ako je potrebno ovdje se moze dodati predgovor (ali to je vise primjereno
%dokt.disertacijama)
% \begin{predgovor}
% Ovdje \'{c}e se opisati \v{s}to me i za\v{s}to navelo na ovaj rad. Mogu\'{c}e je i
% ovdje staviti zahvale.
% 
% Tekst tekst tekst tekst tekst tekst.
% \end{predgovor}

%-- sadrzaj:
\clearpage
\ifpdf
	\phantomsection
\fi
\addcontentsline{toc}{chapter}{{\numberline{}\contentsname}}
\tableofcontents

%
%-- popis slika: ukoliko ih ima
\clearpage
\ifpdf
	\phantomsection
\fi
\addcontentsline{toc}{chapter}{{\numberline{}\listfigurename}}
\listoffigures

%-- popis tablica: ukoliko ih ima
\clearpage
\addcontentsline{toc}{chapter}{{\numberline{}\listtablename}}
\listoftables

%-- ako se koristi klasicni popis oznaka
%\addcontentsline{toc}{chapter}{{\numberline {}\popoznime}}
%	\markboth{\MakeUppercase\popoznime}{}
%\input popozn.tex

\clearpage
\ifpdf
	\phantomsection
\fi
\addcontentsline{toc}{chapter}{{\numberline {}\popoznime}}
\markboth{\MakeUppercase\nomname}{\MakeUppercase\nomname}
\printnomenclature

\clearpage
\addtocounter{brojstrana}{\arabic{page}}\addtocounter{brojstrana}{-1}
    % nuzna komanda za odredjivanje ukupnog broja str.(rimskih i arapskih)
\pagenumbering{arabic} % mjenjam brojke stranice: ali od 1

%
\begin{sazetak}
\noindent Kratki sa\v{z}etak rada na hrvatskom jeziku: najvi\v{s}e jedna stranica
\ldots Tekst tekst tekst tekst tekst tekst.
\end{sazetak}

%
\begin{abstract}
\noindent Short summary of the thesis in one foreign language (english or
german or french): up to one page \ldots Tekst tekst tekst tekst
tekst tekst.
\end{abstract}


%-- ovdje pocinje pravi tekst: najelegantnije jest ubaciti ga kao posebne .tex
%datoteke - jednu za cijeli tekst ili po jednu za svako poglavlje ili ...
\clearpage
\input uvod.tex
\addtocounter{brojslika}{\arabic{figure}}
\addtocounter{brojtablica}{\arabic{table}}

%\input teorija.tex
%\addtocounter{brojslika}{\arabic{figure}}
%\addtocounter{brojtablica}{\arabic{table}}
% ... i tako dalje za ostala poglavlja


\input zakljuc.tex
\addtocounter{brojslika}{\arabic{figure}}
\addtocounter{brojtablica}{\arabic{table}}

\appendix
\chapter{Moj prvi dodatak}
\addtocounter{brojslika}{\arabic{figure}}
\addtocounter{brojtablica}{\arabic{table}}
Tu dolazi prvi prilog odnosno dodatak tekstu (slike, podaci, kod, detaljni
opisi, tablice, \ldots
\newpage 
Pa ni na ovoj stranici!
\section{Malo poglavlje malog dodatka}
\subsection{i još manje podpoglavlje}
Ovo poglavlje poslu\v{z}it \'{c}e za uvod u problem koji se rje\v{s}ava u ovom
radu, te se postavlja odgovaraju\'{c}a hipoteza za doktorsku
disertaciju.

\section{Primjer potpoglavlja}
Tekst tekst tekst tekst tekst tekst primjer reference
\cite{mastersthesis-minimal}, i jo\v{s} jedan \cite{inbook-full}.
Tekst tekst tekst tekst tekst tekst 

Malo ću se pozvati (\eqref{eq:ls-a1}) kao u šđč枊ĐČĆŽ na 
\begin{equation}\label{eq:ls-a1}
	C_l(y)=\frac{2\Gamma(y)}{\Vinf\;c(y)}\:,
\end{equation}
\nomenclature[acal]{$C_l$}{lokalni koeficijent uzgona}%\refeq}%
\nomenclature[ecax]{$c$}{duljina tetive, [m]}%

\subsection{Primjer potpotpoglavlja}

Slijedi prvi primjer slike: FSB.(pogl.sliku~\ref{figa1})
\nomenclature[Kf]{$FSB$}{Fakultet strojarstva i brodogradnje}%

\begin{figure}[h]
  \centering
  \includegraphics[height=1.2cm]{fsb_logo}\\
  \hangcaption{Primjer slike: kod ove slike primjenjeno je
  \emph{vise\'{c}e} zaglavlje -- hangcaption}
  \label{figa1}
\end{figure}

\clearpage
Slijedi mali primjer slike: SUZ.(pogledaj sliku~\ref{figa2})
\nomenclature[Ks]{$SUZ$}{Sveučilište u Zagrebu}%
%prvim slovom unutar uglate zagrade 'K' definiram da li je kratica ili indeks ili
%akcent, a drugim slovom 's' mogu regulirati poredak pojavljivanja u popisu
%oznaka (po abecedi, naravno)

\begin{figure}
  \centering
  \includegraphics[height=2.2cm]{suz_grb}\\
  \hangcaption{Primjer slike: kod ove slike primjenjeno je
  \emph{vise\'{c}e} zaglavlje -- hangcaption}
  \label{figa2}
\end{figure}


Pored slike dan je i primjer tablice (\ref{tablicaa}).

\begin{table}[!h]
\hangcaption{Primjer tablice}
\label{tablicaa}
  \centering
\begin{tabular}{|c c|} \hline
    A & 1 \\
    B & 2 \\
    C & 3 \\
    D & 4 \\ \hline
\end{tabular}
\end{table}

\chapter{i drugi \dots}
\addtocounter{brojslika}{\arabic{figure}}
\addtocounter{brojtablica}{\arabic{table}}
\dots Opet ništa osim ove slike:~\ref{dodatna}
\begin{figure}
  \centering
  \includegraphics[height=1.2cm]{fsb_logo}
  \hangcaption{Primjer slike u prilogu}
  \label{dodatna}
\end{figure}
\newpage
\section{Samo prašim po testu}
Ovo poglavlje poslu\v{z}it \'{c}e za uvod u problem koji se rje\v{s}ava u ovom
radu, te se postavlja odgovaraju\'{c}a hipoteza za doktorsku
disertaciju.


\nomenclature[Kf]{$dod$}{dummy dodatak}%

% literatura:
\newpage
\ifpdf
	\phantomsection
\fi
\bibliographystyle{unsrt} %abbrv
\addcontentsline{toc}{chapter}{{\numberline{}\bibname}}
\bibliography{xampl}

\end{document}               % End of document.

% makeindex glavni.nlo -s nomencl.ist -o glavni.nls 


